\chapter*{Conclusion générale}
\phantomsection
\addcontentsline{toc}{chapter}{Conclusion générale}
\markboth{Conclusion générale}{Conclusion générale}

Ce mémoire avait pour objectif de concevoir et de réaliser une solution bancaire
intelligente centrée sur deux volets complémentaires : la gestion de la relation
client et les services d'investissement. À travers la plateforme \textbf{I-SIB},
le travail mené a permis de proposer une base applicative modulaire, sécurisée
et évolutive, adaptée à un contexte bancaire où la centralisation des
informations et la continuité des parcours deviennent essentielles.

L'analyse des besoins et des acteurs du système a constitué le point de départ,
d'où ont découlé les exigences fonctionnelles et non fonctionnelles.
L'architecture microservices était fixée comme cadre du projet ; nous l'avons
organisée en services métier distincts, placés derrière une passerelle API. Pour
la sécurité et la traçabilité, la plateforme s'appuie sur une gestion des accès
par rôles et sur un service d'audit. Enfin, pour garder une étude lisible, nous
avons resserré la conception détaillée sur deux scénarios, la création d'un
prospect et le passage d'ordres.

Les résultats obtenus couvrent le périmètre fonctionnel réalisé. Le module
Relation Client permet de centraliser les informations utiles au suivi
commercial, de gérer les prospects et d'exploiter des indicateurs d'aide au
suivi, notamment autour du scoring et du risque de départ. Le module SGI permet
la consultation des portefeuilles, des instruments BRVM \cite{brvm_actions} et des informations
nécessaires à la préparation des opérations d'investissement. La plateforme a
également été déployée dans un environnement de validation sur un VPS
AlmaLinux \cite{almalinux_official}, avec une organisation basée sur Docker \cite{docker_docs} et
Docker Compose \cite{docker_compose_docs} afin de tester progressivement les
services réalisés.

Le travail réalisé reste toutefois à consolider avant une exploitation en
production. Les cotations BRVM \cite{brvm_actions} proviennent encore d'une
source publique tierce \cite{brvm_data_public}, sans branchement direct sur le
flux de marché officiel ; les modèles d'intelligence artificielle devront être évalués sur des
données réelles, et certaines connexions avec les autres modules bancaires
restent à finaliser. Le passage en production nécessitera également des tests
complémentaires, un audit de sécurité, une supervision renforcée, des
sauvegardes et une stratégie de reprise clairement définie.

En définitive, ce projet montre la faisabilité d'une plateforme bancaire
unifiée associant relation client, services d'investissement, architecture
modulaire et premiers mécanismes d'aide intelligente. I-SIB constitue ainsi une
base de travail solide, appelée à évoluer progressivement vers une solution plus
complète, mieux intégrée aux données réelles et adaptée aux besoins opérationnels
des institutions financières.
